\expandafter\def\csname HistoricalNoteText349\endcsname{Es wird hier stillschweigend der Fall vorausgesetzt, wo die passive Nadel anfangs in Ruhe oder sehr gelinder Bewegung ist; träfe aber der Anfang der Induction die passive Nadel schon schwingend an, so kann, nach Maassgabe der Stellen der Schwingungsperioden beider Nadeln zur Zeit des Anfangs, auch zuerst eine stetige Abnahme des Schwingungsbogens eintreten, und erst wenn dieser dadurch gleichsam absorbirt ist, wird dann die stetige Zunahme erfolgen.}
\par\clearpage
\SourceLeaf{342}\GaussVSourceFolio{336}{L}{ERDMAGNETISMUS UND MAGNETOMETER.}
\noindent selbst ganz kleine mit allen ihren in den kürzesten Zeitfristen wechselnden Nuancen, selbst an weit von einander. entlegenen Orten, eine ganz bewundernswürdige. Harmonie zeigen. Es sind davon schon mehrere Proben in graphischen Darstellungen bekannt gemacht, von welchen wir hier nur die am 5. und 6. November in Copenhagen und Mailand während 44 Stunden ununterbrochen verfolgten Beobachtungen, und die in zwei Abendstunden des 1. April 1835 in Copenhagen, Altona, Göttingen, Leipzig und Rom angestellten erwähnen wollen. Diesem Vereine zu magnetischen Beobachtungen, an jährlich sechs im Voraus festgesetzten Terminen, schliessen sich schon immer mehr Theilnehmer an; binnen Jahresfrist wird er schon in den entferntesten Theilen des russischen Reichs Mitarbeiter haben. Es steht zu erwarten, dass solche vereinte Bestrebungen uns in Zukunft nähere Aufschlüsse über die räthselhaften Kräfte geben werden, deren Wirkungen sich in gleichem Augenblick über den halben Durchschnitt von Europa verbreiten.

Wir haben hier von dem reichen Stoff, welchen die an dem Magnetometer zu beobachtenden und zu messenden reinmagnetischen Erscheinungen darbieten, nur Einiges ausheben können; jener Apparat ist aber zugleich ein eben so nützliches Werkzeug für die electro-magnetischen Phänomene. Die glänzenden Entdeckungen \textsc{Oersted}'s und \textsc{Faraday}'s haben der Naturforschung eine neue Welt geöffnet, deren Zaubergärten uns mit Bewunderung erfüllen; unterwürfig machen können wir uns diese reichen Gebiete nur unter Führung der Messkunst.

Das erste Erforderniss ist ein Mittel, die Stärke eines galvanischen Stroms durch seine electro-magnetische Wirkung mit Leichtigkeit und Schärfe zu messen. Man bedient sich dazu einer Vorrichtung, die man einen Multiplicator nennt. Es ist dies ein Metalldraht, der in zahlreichen Umwindungen um einen vierseitigen Rahmen geführt, und in die galvanische Kette so gebracht ist, dass er selbst einen Theil des Leitungsdrahts ausmacht. Die einzelnen Windungen dürfen einander nicht metallisch berühren, was man gewöhnlich dadurch verhütet, dass man zu dem Multiplicatordraht solchen anwendet, der mit Seide übersponnen ist. Die einzelnen Umwindungen können hier als unter sich parallele Vierecke betrachtet werden; beim Gebrauch ist der Multiplicator so gestellt, dass die Ebene dieser Vierecke vertical, und nach der verschiedenen Anwendungsart entweder im magnetischen Meridian, oder rechtwinklig dagegen steht. Im innern offenen Raume des Multiplicators befindet sich eine an einem Faden frei
\par\clearpage
\SourceLeaf{343}\GaussVSourceFolio{337}{R}{ERDMAGNETISMUS UND MAGNETOMETER.}
\noindent schwebende Magnetnadel, deren Stellung oder Bewegung bloss vom Erdmagnetismus geregelt wird, so lange den Multiplicatordraht noch kein galvanischer Strom durchläuft. Sobald aber die Kette geschlossen ist, übt der den Multiplicatordraht durchlaufende Strom auf die Nadel eine Kraft aus, deren Richtung immer rechtwinklig gegen die Fläche des Multiplicators ist; der Sinn dieser Richtung ist aber in Beziehung auf die beiden Pole der Nadel entgegengesetzt. In dem erstern der beiden vorhin unterschiedenen Fälle zeigt sich daher die Wirkung der Kraft in einer Ablenkung der Nadel vom magnetischen Meridian, deren Grösse als Maass der Stärke des Stroms betrachtet werden kann, wenigstens davon abhängt. Man hat bisher immer nur äusserst leichte Nadeln angewandt, wobei man zwar grosse Ablenkungen erhielt, die jedoch auf dem angebrachten eingetheilten Kreise sich nur gröblich messen liessen. Eine etwas grössere Genauigkeit kann man durch die zweite von \textsc{Fechner} angewandte Einrichtung erhalten, wo die Richtung der von dem galvanischen Strom ausgeübten Kraft in dem magnetischen Meridian selbst liegt, und folglich ihre Wirkung (je nach der Richtung des Stromes im Draht) entweder den Erdmagnetismus verstärkt, oder verringert, und wodurch also die Schwingungsdauer der Nadel entweder kürzer oder länger wird, als sie unter dem reinen Einfluss des Erdmagnetismus war. Dieses Verfahren hat indess, abgesehen davon, dass die Schärfe noch immer lange nicht so gross ist, als man wünschen muss, das Unangenehme, dass es, da man zur Bestimmung der Schwingungsdauer eine beträchtliche Anzahl von Schwingungen zu beobachten hat, sehr mühsam wird, und, insofern die Stromstärke während der Dauer eines Versuches veränderlich ist, nur eine Art Mittelwerth angibt; zur Messung der Stärke solcher Ströme, die, wie die durch die sogenannte Induction hervorgebrachten (von denen später noch die Rede sein wird), nur wenige Augenblicke dauern, ist diese Methode gar nicht anzuwenden.

Man kann nun aber leicht das Magnetometer zu einem eben so bequemen, als scharfen Galvanometer einrichten, wenn man es mit einem Multiplicator verbindet, dessen Ebene, wie bei der ersten vorhin erwähnten Einrichtung, im magnetischen Meridian ist. Da in dem Magnetometer immer grosse Stäbe angewandt werden, so ist sowohl die innere Weite des Multiplicators, als seine Drahtlänge viel grösser, als bei den sonst gebräuchlichen Multiplicatoren. Der erstere Umstand trägt dazu bei, die Einwirkung des galvanischen Stroms auf die Nadel, der andere hingegen, die Intensität des galvanischen Stroms selbst schwächer zu
\par\GaussVSourceGathering{43}
\par\clearpage
\SourceLeaf{344}\GaussVSourceFolio{338}{L}{ERDMAGNETISMUS UND MAGNETOMETER.}
\noindent machen; aus beiden Ursachen finden daher im Allgemeinen keine so grosse Ablenkungen der Magnetnadel Statt, wie bei den andern Galvanometern. Dagegen aber kann man hier die Grösse der mässigen Ablenkung mit äusserster Schärfe messen. Das Magnetometer der Göttinger Sternwarte hat einen Multiplicator von 270, das des magnetischen Observatorium einen von 200 Umwindungen; die Drahtlänge des erstern ist 2700, die des andern 1100 Fuss; beide hängen durch eine 450 Fuss lange Drahtverbindung unter sich, und durch eine 6000 Fuss lange, noch mit einem Paar ähnlicher, obwohl etwas kleinerer Apparate, in dem eine Viertelstunde davon entfernten physikalischen Kabinet zusammen, so dass ein galvanischer durch diese ganze bisher in ihrer Art einzige Kette getriebener Strom eine Drahtlänge von fast einer halben Meile zu durchlaufen hat. Und doch bewegt ein solcher Strom, nur von einem kleinen Plattenpaar mit blossem Brunnenwasser erregt, in allen vier Apparaten die Magnetnadeln zu Ausschlägen von vielen hundert Scalentheilen; der Strom durchläuft diese Strecke in einer ganz unmessbar kleinen Zeit, so dass durch Beobachtung des Anfangs der Bewegung der Magnetnadeln die Uhren an den vier Plätzen schärfer als durch irgend ein anderes Mittel mit einander verglichen werden können. Durch eine Vorrichtung. die man einen Commutator oder Gyrotrop nennt, kann man die Richtung des Stroms augenblicklich in die entgegengesetzte verwandeln, oder auch den Strom selbst unterbrechen, was dann auf die Bewegung der Nadeln einen entsprechenden Einfluss hat. Man ist durch diese Vorrichtungen über die Bewegungen so sehr Herr, dass man sich ihrer zu telegraphischen Zeichen bedienen kann, die ganz unabhängig von Tageszeit und Witterung in verschlossenem Zimmer gegeben, und eben so empfangen werden. Öftere Versuche, ganze Wörter und kleine Phrasen auf diese Weise zu signalisiren, haben den vollkommensten Erfolg gehabt. Was hier nur ein interessanter physicalischer Versuch ist, liesse sich, wie man mit Zuversicht voraussagen kann, bei einer Ausführung in noch viel grösserem Maassstabe, und unter Anwendung starker galvanischer Säulen oder sonstiger electrometrischer Kräfte, starker Multiplicatoren und starker Leitungsdrähte zu telegraphischen Verbindungen auf zehn, zwanzig und mehrere Meilen in einem Schlage, benutzen. Es ist Hoffnung, dass schon in Kurzem ein ähnlicher Versuch auf mehrere Meilen Entfernung durch einen eifrigen und kenntnissvollen Freund der Naturwissenschaften ausgeführt werden wird. Könnte man, unbeschadet anderer zu nehmender Rücksichten, die einzelnen Schienen der Eisenbahnen sicher
\par\clearpage
\SourceLeaf{345}\GaussVSourceFolio{339}{R}{ERDMAGNETISMUS UND MAGNETOMETER.}
\noindent und leicht metallisch verbinden, so würden diese mit Vortheil anstatt der Leitungsdrähte dienen können. Überhaupt scheint einer Erstreckung der electro-magnetischen Telegraphie, selbst auf ungeheure Entfernungen, nichts im Wege zu stehen, als der Anwachs der Kosten, da grössere von dem galvanischen Strom ohne Zwischenstation zu durchlaufende Strecken zugleich dickere Leitungsdrähte erfordern.

Wir haben oben \textsc{Faraday} neben \textsc{Oersted} genannt; beider Entdeckungen haben in der Naturwissenschaft Epoche gemacht; sie sind auf das engste mit einander verbunden, ja die eine ist, wie an einem andern Orte näher nachgewiesen werden soll, als das vollkommene Seitenstück der andern zu betrachten. \textsc{Oersted} entdeckte die Einwirkung eines schon bestehenden galvanischen Stromes auf die magnetischen Stoffe; \textsc{Faraday} fand, dass, indem die magnetischen Stoffe sich neben einem zur Leitung eines galvanischen Stromes fähigen Körper bewegen, in diesem ein solcher Strom hervorgebracht wird, der aber nur so lange dauert, wie eben jene Bewegung der magnetischen Stoffe. Ohne in die genauern Bedingungen hier einzugehen, wollen wir nur bemerken, dass gleiche Bewegungen der beiden entgegengesetzten magnetischen Flüssigkeiten entgegengesetzte galvanische Ströme erzeugen, also ihre Wirkungen sich selbst neutralisiren, wenn jene gleichzeitig sind. Daher bringt die Bewegung eines Trägers der magnetischen Flüssigkeiten, in welchem sie noch nicht geschieden sind, des Eisens oder des nicht magnetisirten Stahls, keinen galvanischen Strom im benachbarten Metall hervor, wohl aber der Act der Scheidung selbst, wenn z.~B. weiches Eisen durch plötzliches Anfügen an die Pole eines Hufeisenmagnets, oder durch irgend ein anderes Mittel plötzlich magnetisch gemacht wird; und eben so muss wieder das plötzliche Abreissen, nach welchem die im Eisen getrennt gewesenen magnetischen Flüssigkeiten sich wieder vereinigen, einen galvanischen Strom von der der vorigen entgegengesetzten Richtung hervorbringen. Die auf diese Weise erzeugten galvanischen Ströme sind (wie der Act der Scheidung oder Wiedervereinigung der magnetischen Flüssigkeiten selbst) von äusserst kurzer Dauer, aber, wenn man die übrigen Umstände zweckmässig anordnet, von grosser Intensität, so dass man dadurch Funken und andere mit starken galvanischen Strömen verbundene Erscheinungen hervorgebracht hat, welche das Erstaunen der Liebhaber der Physik erregen. Eine andere Art, den magnetischen Flüssigkeiten ungleiche Bewegungen zu ertheilen (was immer die Bedingung dieser Stromerregungsart bleibt),
\par\GaussVSourceGathering{43\textsuperscript{*}}
\par\clearpage
\SourceLeaf{346}\GaussVSourceFolio{340}{L}{ERDMAGNETISMUS UND MAGNETOMETER.}
\noindent besteht aber darin, dass man solche Träger derselben, in welchem sie schon geschieden sind (einen Magnetstab, oder eine Verbindung mehrerer), entweder selbst auf eine zweckmässige Art relativ gegen einen nahen Leiter bewegt, oder auch, was in der Wirkung ganz einerlei ist, jene Träger ruhen lässt, und den Leiter, der den Strom empfangen soll, bewegt.

Wesentlich sind diese beiden Arten von Stromerregung (Induction) gar nicht verschieden; die zweite ist aber allein brauchbar für solche Versuche, bei welchen es um genaue Kenntniss der Grössenverhältnisse zu thun ist. Man kann sich dazu eines sehr einfachen Mittels bedienen.

Eben so wie man zur Verstärkung des von \textsc{Oersted} entdeckten Einflusses des galvanischen Stroms auf die Magnetnadel einen zu zahlreichen Windungen geformten Leitungsdraht (Multiplicator) anwendet, verstärkt man den Strom, welchen die relative Ortsveränderung des den Strom empfangenden Drahts gegen den Magnet erzeugt, dadurch, dass viele Theile des Drahts auf gleiche Weise afficirt werden. Eine dazu dienende Vorrichtung kann man einen Inductions-Multiplicator, oder schlechthin einen Inductor nennen. Ein solcher bei dem Apparat der Göttinger Sternwarte gebrauchter Inductor besteht in einer cylindrischen Rolle, im Lichten beinahe vier Zoll weit, um deren äussere Fläche ein mit Seide übersponnener Kupferdraht 3537 mal (in einer Länge von etwa 3600 Fuss) gewunden ist, dessen Enden mit der Kette in Verbindung gebracht sind. Zwei starke Magnetstäbe, jeder von 25 Pfund, sind zu Einem kräftigen Magnet verbunden. Das blosse Aufschieben der Rolle auf diesen Magnet bis zu dessen Mitte bewirkt in dem Draht und der ganzen damit verbundenen Kette, mithin auch in den verschiedenen Multiplicatoren, welche Theile davon ausmachen, einen kräftigen galvanischen Strom, welcher also entsprechende Bewegungen in denjenigen Magnetnadeln hervorbringt, welche sich in den betreffenden Multiplicatoren befinden, und dessen Stärke durch die Magnetometer scharf gemessen wird. Der Strom dauert immer nur so lange, wie die Bewegung der Inductionsrolle. Das Wiederabziehen, und eben so das Verkehrt-Wiederaufschieben, bewirkt einen dem vorigen entgegengesetzten Strom; vermittelst der in der Kette befindlichen Commutatoren hat man in seiner Gewalt, dem Strom in den Multiplicatoren jedesmal eine beliebige Richtung zu geben. Es ist hiebei ein höchst wichtiger Umstand, dass, obgleich die Stärke des Stroms von der Geschwindigkeit der Bewegung der Rolle abhängt, dennoch (weil die \textit{Dauer} desto kürzer ist, je schneller man mit
\par\clearpage
\SourceLeaf{347}\GaussVSourceFolio{341}{R}{ERDMAGNETISMUS UND MAGNETOMETER.}
\noindent der Operation zu Ende kommt) die Gesammtwirkung auf die Bewegung der Magnetnadeln in den Multiplicatoren von der Schnelligkeit der Bewegung fast ganz unabhängig bleibt, insofern diese in einer oder ein paar Secunden vollendet wird. Beim Gebrauch lässt man gewöhnlich auf ein Abziehen der Inductionsrolle ein verkehrtes Wiederaufschieben unmittelbar folgen, was zusammen ein Wechsel heissen kann. Die Wirkung eines solchen Wechsels, auch wenn der Strom durch die ganze jetzt fast 15000 Fuss lange Kette getrieben wird, ist so stark, dass die betreffenden Magnetnadeln Bewegungen dadurch erhalten, die viele hundert Scalentheile betragen. Man kann aber in kurzer Zeit sehr viele solche Wechsel eintreten lassen, die vermöge entsprechenden Spiels des Commutators alle einander verstärken, und die Magnetnadeln der Magnetometer in so grosse Bewegungen wie man will, versetzen. Die Erfahrung zeigt bei solchen Versuchen eine Übereinstimmung in den quantitativen Verhältnissen, die nichts zu wünschen übrig lässt, und die Erforschung der Gesetze dieser so höchst interessanten Naturphänomene eben so sehr befestigt als erleichtert hat.

Diese Gesetze, zu deren Entwicklung hier nicht der Ort ist, bestätigen sich überall so vollkommen, dass man den Erfolg von Versuchen, sobald man die Umstände, von welchen sie abhängen, nach ihren Grössenverhältnissen kennt, so sicher im Voraus bestimmen kann, wie die Erscheinungen am Sternenhimmel. Einen solchen Versuch, der zu den auffallendsten im Gebiet des Electromagnetismus gehört, wollen wir hier noch anführen.

Eben so gut, wie durch die relative Bewegung der Inductionsrolle gegen den Hülfsmagnet. in dem Draht der erstern, wenn er eine wo immer geschlossene Kette bildet, ein galvanischer Strom hervorgerufen wird, ist auch, jenen Inductor ganz bei Seite gesetzt, die Schwingungsbewegung einer Magnetnadel in ihrem Multiplicator, sobald dieser eine geschlossene Kette darstellt, oder einen Theil davon ausmacht, von einem galvanischen Strome in dieser begleitet, nur ist dieser den Umständen nach viel schwächer als jener. Betrachten wir z. B. den fünfundzwanzigpfündigen Magnetstab des Magnetometers der Göttinger Sternwarte als schwingend, so ist der durch \textit{seine} Inductionswirkung erzeugte Strom schwächer, als der am Inductor hervorgebrachte, erstlich weil in jenem Fall nur Ein Magnetstab wirksam ist, im andern zwei von derselben Grösse; zweitens weil der jenen Strom empfangende Multiplicatordraht nur 270 Umwindungen bildet, der Inductordraht aber 3537; drittens weil die Windungen des letztern viel
\par\clearpage
\SourceLeaf{348}\GaussVSourceFolio{342}{L}{ERDMAGNETISMUS UND MAGNETOMETER.}
\noindent enger sind, als die des erstern; viertens wegen der äusserst langsamen Schwingungsbewegung der Magnetometernadel, da theils der schmale, den Stab einschliessende Kasten nur Schwingungen von mässiger Grösse verstattet, theils jede Schwingung eine so lange Zeit erfordert. Wie schwach aber auch der erstere Strom, verglichen mit dem andern, ist, so tritt doch seine Existenz sehr bestimmt hervor. Er muss nemlich, gleich jedem andern, wie immer erzeugten, den Multiplicator durchlaufenden Strom, auf die im Multiplicator befindliche Nadel wirken, und diese Rückwirkung zeigt sich, ganz der Theorie gemäss, darin, dass der Schwingungsbogen viel rascher abnimmt, als bloss vermöge des Widerstandes der Luft, oder wenn gar kein Strom da ist, geschehen würde, gleichsam, als schwänge die Nadel in einer vielfach dichtern Flüssigkeit, als die Luft ist. Dies bestätigt die Erfahrung vollkommen. Ja diese allmählige Lähmung der Bewegung (wenn wir uns des Ausdrucks bedienen dürfen) ist am stärksten, wenn die Kette gleich hinter dem Multiplicator abgeschlossen, weniger stark, wenn eine grössere Drahtlänge noch mit in die Kette gebracht ist, am geringsten, wenn die ganze 15000 Fuss betragende Drahtlänge Eine Kette bildet, aber auch dann noch immer sehr beträchtlich; sobald aber die Kette wo immer geöffnet ist, fällt \textit{dieser} Einfluss ganz weg, und die Abnahme der Grösse des Schwingungsbogens reducirt sich sogleich auf den geringen Betrag, der hauptsächlich dem Widerstande der Luft zuzuschreiben ist, und auch dann noch bleibt, wenn man den Multiplicator ganz weggenommen hat. Übrigens ist diese Rückwirkung des galvanischen Stroms auf die Nadel, durch deren Schwingung er selbst erzeugt wird, auch schon bei kleinern Nadeln sehr bestimmt zu bemerken, wenn nur der sie eng umgebende Multiplicator viele Umwindungen, oder auch, wenn nur bei weniger Umwindungen der Draht eine beträchtliche Stärke hat, vorausgesetzt, dass in letzterm Fall die Kette gleich hinter dem Multiplicator abgesperrt ist, ja eigentlich beruhen die ähnlichen Erscheinungen, die, schon vor Entdeckung der Induction, \textsc{Arago} an Magnetnadeln, die über Metallplatten schwingen, bemerkt hat, auf demselben Grunde.

Allein noch viel klarer tritt das Dasein eines auf diese Art erzeugten galvanischen Stromes hervor, wenn, wie bei den Göttinger Einrichtungen, noch andere Magnetometer sich in der verlängerten Kette befinden. Wenn man den grossen Magnetstab des Magnetometers der Sternwarte in etwas beträchtliche Schwingungen versetzt, so nehmen diese, falls die Kette noch nicht geschlossen
\par\clearpage
\SourceLeaf{349}\GaussVSourceFolio{343}{R}{ERDMAGNETISMUS UND MAGNETOMETER.}
\noindent ist, nur sehr langsam an Grösse ab, und die Nadeln der Magnetometer im physikalischen Kabinet und im magnetischen Observatorium bleiben in Ruhe, oder in derjenigen regelmässigen Schwingungsbewegung, die sie eben haben; allein von dem Augenblick an, wo die Kette geschlossen wird, fangen nicht bloss die Schwingungsbögen des grossen Stabs sogleich an, viel schneller abzunehmen, sondern wie durch eine magische Sympathie kommen die andern Nadeln mit in Bewegung, falls sie vorher ganz in Ruhe waren, oder, wenn sie sich schon selbst in Bewegung befanden, erhalten ihre Schwingungen einen andern Charakter, so dass zweierlei Schwingungen sich gleichsam vermengen, die ihnen natürlichen, mit der denselben zukommenden Schwingungsperiode, und die inducirten, ihnen gleichsam aufgedrungenen, die eine ganz andere Periode haben. Die Periode dieser inducirten Schwingungen ist an Dauer der Schwingungsperiode des grossen inducirenden Stabs genau gleich (42 Secunden), allein immer fällt ihr Anfang und Ende der Zeit nach nicht mit Anfang und Ende einer Schwingung oder Rückschwingung des grossen Stabes zusammen, sondern vielmehr mit der Mitte einer solchen. Was endlich die Grösse der vermengten Schwingungen betrifft, so ist die der natürlichen Schwingungen abhängig von dem Bewegungszustande der Nadel beim Anfang der Induction, hingegen die Grösse der inducirten Schwingungen von diesem Initialzustande ganz unabhängig, und bloss durch die Grösse der Schwingungen des inducirenden Stabs bestimmt.

Ein fast noch merkwürdigerer Erfolg findet aber Statt, wenn in einem der andern Apparate eine Nadel aufgehängt ist, deren natürliche Schwingungsdauer auch 42 Secunden beträgt, oder, um es allgemein auszudrücken, der Schwingungsdauer der inducirenden Nadel genau gleich ist. In diesem Fall behalten die sympathetisch inducirten Schwingungsbewegungen dieselbe Periode, aber sie werden immer grösser\HistoricalNote{349}. Schon sehr oft und mit immer gleichem Erfolg ist dieser interessante Versuch angestellt, und das wunderbar erscheinende Schauspiel beobachtet, dass ein Magnetstab lediglich durch einen andern in so grosser Entfernung schwingenden angeregt, aus seiner Ruhe gerissen und zu immer schnel-\par\clearpage
\SourceLeaf{350}\GaussVSourceFolio{344}{L}{ERDMAGNETISMUS UND MAGNETOMETER.}
\noindent lerer Bewegung angespornt wird. Schon nach wenigen Minuten wurde bei diesen Versuchen die Bewegung so gross, dass die Scale des Magnetometers nicht mehr ausreichte, sie unmittelbar zu messen; aber mittelbarer Weise konnte man sich doch leicht, länger als eine Stunde, von der beständig fortdauernden Beschleunigung überzeugen, und in der That muss die Zunahme der Bewegung so lange fortdauern. bis die andern Ursachen, die zur Schwächung der Bewegung wirken, der Vergrösserungsursache das Gleichgewicht halten, von welchem Zeitpunkt an dann die Bögen allmählig wieder kleiner werden.

Wir finden hier im Kleinen eine Art von Spiegelung der gegenseitigen Einwirkung der Himmelskörper; der erste Versuch erinnert uns an die periodischen, der andere an die Säcularstörungen, welche ein Planet an einem andern ausübt. Aber wie diese Störungen in allen ihren Verwicklungen aus dem allgemeinen Gravitationsgesetze als nothwendige Folgen hervorgehen, so folgen auch die hier erzählten Erscheinungen von selbst aus einem allgemeinen sehr einfachen electro-magnetischen Grundgesetze; auch waren sie mit allen begleitenden hier nur in allgemeinen Umrissen angedeuteten Umständen aus diesem Grundgesetze durch die Theorie im Voraus abgeleitet, ehe die Versuche selbst angestellt wurden.

\par\vspace{4.5\baselineskip}\centerline{\rule{.255\linewidth}{.35pt}}\par
\clearpage
\SourceLeaf{351}\GaussVSourceUnnumbered
\vspace*{1.35in}
\begin{center}
{\Large\addfontfeatures{LetterSpace=12}EINLEITUNG}\par\medskip
\resizebox{.96\linewidth}{!}{\shortstack[c]{[FÜR DIE ZEITSCHRIFT: RESULTATE AUS DEN BEOBACHTUNGEN DES MAGNETISCHEN\\[.4ex]VEREINS IM JAHRE 1836. HERAUSGEGEBEN VON C. F. GAUSS UND W. WEBER.]}}
\par\medskip\rule{.09\linewidth}{.35pt}
\end{center}
\vspace{1.2\baselineskip}
Unter den mannigfaltigen Äusserungen der erdmagnetischen Kraft, deren Ergründung nur durch zahlreiche an den verschiedensten Punkten der Erdoberfläche fortgesetzt anzustellende genaue Beobachtungen zu erreichen ist, bedürfen die \textit{unregelmässigen Änderungen}, welchen wir jene Kraft unterworfen finden, am meisten eines streng geordneten Zusammenwirkens der Beobachter. Es ist bekannt genug, dass die Bestimmungsstücke der erdmagnetischen Kraft, die Abweichung, die Neigung und ohne Zweifel auch die Stärke (wenn gleich in Beziehung auf die letzte, die erst seit einigen Jahrzehnden in den Kreis der Forschungen aufgenommen ist, noch hinlängliche Erfahrungen fehlen) fortwährend Veränderungen erleiden, seculäre erst nach längerer Zwischenzeit in die Augen fallende, aber im Laufe der Zeit sehr beträchtlich werdende, und periodische nach den Jahres- und Tageszeiten wechselnde. Aber so weit diese Veränderungen mit Regelmässigkeit geschehen, ist ein streng geordnetes Zusammenwirken der Beobachter an verschiedenen Orten, wenn auch für die Beschleunigung der Erweiterung unserer Einsicht höchst wünschenswerth, doch nicht wesentlich nothwendig, und jeder Beobachter kann auch unabhängig von den andern nützliche Beiträge liefern.

Anders verhält es sich hingegen mit den unregelmässigen Veränderungen, welchen man erst in den letzten Jahren eine grössere Aufmerksamkeit zu widmen angefangen hat. Dass während der Sichtbarkeit eines Nordlichts die Magnetnadel unregelmässige und oft sehr grosse Bewegungen zeigt, haben schon vor bei-
\par\GaussVSourceGathering{44}

\par\clearpage
